\documentclass[11pt,a4paper]{article}

% ============================================================
% PAQUETES
% ============================================================
\usepackage[utf8]{inputenc}
\usepackage[T1]{fontenc}
\usepackage[spanish,es-nodecimaldot]{babel}
\usepackage{lmodern}
\usepackage{geometry}
\usepackage{microtype}
\usepackage{xcolor}
\usepackage{booktabs}
\usepackage{array}
\usepackage{tabularx}
\usepackage{enumitem}
\usepackage{fancyhdr}
\usepackage{titlesec}
\usepackage{hyperref}
\usepackage{graphicx}
\usepackage{parskip}

% ============================================================
% CONFIGURACIÓN
% ============================================================
\geometry{
    top=2.2cm,
    bottom=2.2cm,
    left=2.4cm,
    right=2.4cm
}

\definecolor{orbitblue}{HTML}{17365D}
\definecolor{orbitlight}{HTML}{EAF0F7}
\definecolor{orbitgray}{HTML}{555555}

\hypersetup{
    colorlinks=true,
    linkcolor=orbitblue,
    urlcolor=orbitblue,
    pdftitle={OrbitPay OO -- Executive Summary},
    pdfauthor={Luis Angel Leyva Castillo},
    pdfsubject={Proyecto integrador de Programación Orientada a Objetos}
}

\pagestyle{fancy}
\fancyhf{}
\lhead{\textsc{OrbitPay OO}}
\rhead{\textsc{Executive Summary}}
\cfoot{\thepage}
\setlength{\headheight}{14pt}

\titleformat{\section}
    {\Large\bfseries\color{orbitblue}}
    {\thesection.}
    {0.6em}
    {}

\titleformat{\subsection}
    {\large\bfseries\color{orbitblue}}
    {\thesubsection.}
    {0.5em}
    {}

\setlist[itemize]{
    leftmargin=1.2cm,
    itemsep=2pt,
    topsep=4pt
}

% ============================================================
% DOCUMENTO
% ============================================================
\begin{document}

% ============================================================
% PORTADA
% ============================================================
\begin{titlepage}
    \centering

    \vspace*{2.0cm}

    {\Huge\bfseries\color{orbitblue} OrbitPay OO\par}

    \vspace{0.5cm}

    {\Large Motor de pagos y suscripciones orientado a objetos\par}

    \vspace{1.5cm}

    \rule{\textwidth}{1pt}

    \vspace{0.6cm}

    {\LARGE\bfseries Executive Summary\par}

    \vspace{0.6cm}

    \rule{\textwidth}{1pt}

    \vspace{2cm}

    \begin{tabular}{rl}
        \textbf{Autor:} & Luis Angel Leyva Castillo \\[0.3cm]
        \textbf{Versión:} & 1.0.0 \\[0.3cm]
        \textbf{Tecnología:} & Python 3.12.14 \\[0.3cm]
        \textbf{Licencia:} & MIT \\[0.3cm]
        \textbf{Tipo de entrega:} &
        Proyecto integrador -- Módulo 2
    \end{tabular}

    \vfill

    \begin{minipage}{0.82\textwidth}
        \centering
        \small
        Documento ejecutivo de síntesis de la solución,
        arquitectura, decisiones técnicas, pruebas,
        calidad y release final de OrbitPay OO.
    \end{minipage}

    \vspace{1cm}

    {\small \today}

\end{titlepage}

% ============================================================
% RESUMEN EJECUTIVO
% ============================================================
\section{Resumen ejecutivo}

OrbitPay OO es un motor de pagos y suscripciones desarrollado
en Python mediante Programación Orientada a Objetos, diseñado
como proyecto integrador del Módulo 2.

Su propósito es transformar un procesamiento monolítico de pagos,
caracterizado por lógica condicional anidada, riesgo de cargos
duplicados, inconsistencias de saldo y ausencia de pruebas
automatizadas, en una solución modular, extensible y verificable.

La solución implementa un modelo de dominio compuesto por
cuentas, transacciones y suscripciones, junto con una abstracción
polimórfica para métodos de pago. La arquitectura incorpora los
patrones Factory, Strategy y Observer, así como principios SOLID
y mecanismos de validación, encapsulamiento e idempotencia.

El release final corresponde a la versión \textbf{1.0.0},
distribuible como paquete Python mediante archivos
\texttt{wheel} y \texttt{tar.gz}. El proyecto alcanzó
\textbf{98\,\% de cobertura de código}, cuenta con pruebas
unitarias y de integración, validación estática mediante
\texttt{mypy --strict}, análisis mediante Ruff, formateo
mediante Black y automatización mediante pre-commit y
GitHub Actions.

% ============================================================
% PROBLEMA
% ============================================================
\section{Problema y objetivo}

El escenario inicial de OrbitPay presentaba características
propias de un sistema monolítico: lógica de negocio concentrada,
decisiones basadas en tipos concretos de métodos de pago,
ausencia de separación clara de responsabilidades y falta de
pruebas automatizadas.

Estas características dificultaban la incorporación de nuevos
métodos de pago y aumentaban el riesgo de errores relacionados
con el procesamiento de operaciones y la consistencia del saldo.

El objetivo fue desarrollar una solución orientada a objetos
que permitiera:

\begin{itemize}
    \item encapsular el estado crítico del dominio;
    \item validar las reglas fundamentales del negocio;
    \item utilizar polimorfismo para los métodos de pago;
    \item desacoplar el cálculo de comisiones;
    \item implementar notificaciones mediante eventos;
    \item evitar el procesamiento duplicado de transacciones;
    \item aplicar principios SOLID;
    \item automatizar pruebas y controles de calidad;
    \item generar un paquete Python instalable y reproducible.
\end{itemize}

% ============================================================
% ARQUITECTURA
% ============================================================
\section{Solución y arquitectura}

OrbitPay se estructura en cuatro áreas principales:

\begin{description}
    \item[\textbf{Dominio:}]
    contiene \texttt{Cuenta}, \texttt{Transaccion},
    \texttt{Suscripcion} y las excepciones del negocio.

    \item[\textbf{Pagos:}]
    define la abstracción \texttt{MetodoPago} y sus
    implementaciones \texttt{Tarjeta},
    \texttt{Transferencia} y \texttt{Wallet}.

    \item[\textbf{Patrones:}]
    concentra las implementaciones de Factory,
    Strategy y Observer.

    \item[\textbf{Motor:}]
    \texttt{PagoEngine} coordina el procesamiento de las
    operaciones mediante abstracciones y dependencias
    inyectadas.
\end{description}

El motor no necesita conocer los detalles concretos de cada
método de pago. Recibe una instancia de \texttt{MetodoPago},
una estrategia de comisión y un gestor de eventos, reduciendo
el acoplamiento entre componentes.

% ============================================================
% DIAGRAMA DE ARQUITECTURA
% ============================================================
\section{Diagrama de clases}

La Figura~\ref{fig:class-diagram} presenta la estructura
principal de OrbitPay OO y las relaciones entre las entidades
del dominio, los métodos de pago, las estrategias de comisión,
el motor de pagos y el sistema de eventos.
\begin{figure}[htbp]
    \centering
    \includegraphics[
        width=\textwidth,
        height=0.62\textheight,
        keepaspectratio
    ]{docs/class-diagram.pdf}
    \caption{Diagrama de clases de OrbitPay OO.}
    \label{fig:class-diagram}
\end{figure}

% ============================================================
% POO
% ============================================================
\section{Programación Orientada a Objetos}

El proyecto evidencia los principales conceptos estudiados
durante el módulo.

\subsection{Encapsulamiento}

El saldo de \texttt{Cuenta} se mantiene como estado interno y
se modifica mediante operaciones controladas como
\texttt{depositar()} y \texttt{retirar()}.

El estado de \texttt{Transaccion} también se controla mediante
operaciones de dominio como \texttt{aprobar()} y
\texttt{rechazar()}.

\subsection{Abstracción}

\texttt{MetodoPago}, \texttt{EstrategiaComision} y
\texttt{ObservadorPago} definen contratos mediante clases
abstractas.

\subsection{Herencia}

Los diferentes métodos de pago implementan el contrato común
de \texttt{MetodoPago}, mientras que las estrategias concretas
implementan \texttt{EstrategiaComision}.

\subsection{Polimorfismo}

El motor procesa cualquier implementación válida de
\texttt{MetodoPago} sin utilizar condicionales basados en el
tipo concreto.

\subsection{Composición}

\texttt{PagoEngine} recibe sus dependencias mediante el
constructor, permitiendo sustituir estrategias y gestores de
eventos sin modificar el motor principal.

% ============================================================
% PATRONES
% ============================================================
\section{Patrones de diseño}

\subsection{Factory Method}

\texttt{MetodoPagoFactory} centraliza la creación de métodos
de pago y permite registrar nuevas implementaciones mediante
un registro extensible.

\subsection{Strategy}

\texttt{EstrategiaComision} separa el algoritmo de cálculo de
comisión del motor de pagos.

Se implementaron:

\begin{itemize}
    \item \texttt{ComisionFija};
    \item \texttt{ComisionPorcentual}.
\end{itemize}

\subsection{Observer}

\texttt{GestorEventosPago} permite notificar eventos de pago
sin acoplar directamente el motor con los componentes
consumidores.

Se contemplan los eventos:

\begin{itemize}
    \item \texttt{PagoAprobado};
    \item \texttt{PagoRechazado}.
\end{itemize}

% ============================================================
% SOLID
% ============================================================
\section{Principios SOLID}

\begin{description}
    \item[\textbf{SRP -- Single Responsibility Principle:}]
    responsabilidades separadas entre dominio, métodos de
    pago, cálculo de comisiones, eventos y procesamiento.

    \item[\textbf{OCP -- Open/Closed Principle:}]
    Factory y Strategy permiten extender comportamientos sin
    modificar el motor principal.

    \item[\textbf{LSP -- Liskov Substitution Principle:}]
    las implementaciones de \texttt{MetodoPago} cumplen el
    contrato definido por la abstracción.

    \item[\textbf{ISP -- Interface Segregation Principle:}]
    las interfaces se mantienen pequeñas y enfocadas.

    \item[\textbf{DIP -- Dependency Inversion Principle:}]
    \texttt{PagoEngine} depende de abstracciones y recibe
    sus dependencias mediante inyección.
\end{description}

% ============================================================
% REFACTOR
% ============================================================
\section{Refactorización}

Durante el desarrollo se identificaron y corrigieron tres
problemas concretos.

\subsection{Refactor 1 -- Estado de transacción}

Se sustituyó la modificación directa del estado por atributos
internos y operaciones de dominio
(\texttt{aprobar()} y \texttt{rechazar()}).

\subsection{Refactor 2 -- Factory extensible}

La creación de métodos de pago se transformó en un registro
extensible, permitiendo incorporar nuevos tipos sin modificar
la lógica central de la Factory.

\subsection{Refactor 3 -- Lógica de rechazo duplicada}

La gestión de rechazos se centralizó mediante
\texttt{\_rechazar\_pago()}, eliminando duplicación dentro del
motor.

% ============================================================
% PRUEBAS
% ============================================================
\section{Pruebas y calidad}

El proyecto cuenta con pruebas unitarias para dominio, métodos
de pago y patrones, además de pruebas de integración para
\texttt{PagoEngine}.

La ejecución final alcanzó:

\begin{center}
    {\Large\bfseries 98\,\% de cobertura}
\end{center}

superando el mínimo requerido de 80\,\%.

Los controles automatizados incluyen:

\begin{itemize}
    \item \texttt{pytest};
    \item \texttt{pytest-cov};
    \item Black;
    \item Ruff;
    \item Mypy en modo \texttt{strict};
    \item pre-commit;
    \item GitHub Actions.
\end{itemize}

El pipeline de integración continua ejecuta las verificaciones
de formato, análisis estático, tipado y pruebas con un requisito
mínimo de 80\,\% de cobertura.

% ============================================================
% IDEMPOTENCIA
% ============================================================
\section{Idempotencia y reglas de negocio}

El motor mantiene un registro de transacciones procesadas para
impedir que una misma transacción sea procesada nuevamente.

También se aplican reglas fundamentales del dominio:

\begin{itemize}
    \item los montos deben ser mayores que cero;
    \item el saldo no puede ser negativo;
    \item una cuenta no puede retirar más fondos de los disponibles;
    \item los identificadores obligatorios deben ser válidos;
    \item los estados de las transacciones están restringidos a valores definidos.
\end{itemize}

% ============================================================
% RELEASE
% ============================================================
\section{Empaquetado y release}

OrbitPay fue empaquetado como versión \textbf{1.0.0} mediante
\texttt{python -m build}.

Se generaron los siguientes artefactos:

\begin{itemize}
    \item \texttt{orbitpay-1.0.0-py3-none-any.whl};
    \item \texttt{orbitpay-1.0.0.tar.gz}.
\end{itemize}

La instalación del wheel fue validada en un entorno virtual
independiente del entorno de desarrollo.

La prueba confirmó:

\begin{itemize}
    \item instalación en \texttt{site-packages};
    \item versión instalada \texttt{1.0.0};
    \item importación correcta del paquete;
    \item ejecución funcional de componentes del dominio.
\end{itemize}

% ============================================================
% RESULTADOS
% ============================================================
\section{Resultados}

\begin{center}
\renewcommand{\arraystretch}{1.3}
\begin{tabularx}{0.85\textwidth}{
    >{\bfseries}l
    X
}
\toprule
Indicador & Resultado \\
\midrule
Versión final & 1.0.0 \\
Cobertura & 98\,\% \\
Mínimo requerido & 80\,\% \\
Patrones implementados & 3 \\
Refactors documentados & 3 \\
Tipado & Mypy \texttt{strict} \\
Formateo & Black \\
Análisis estático & Ruff \\
Automatización & Pre-commit + GitHub Actions \\
Distribución & Wheel + Source Distribution \\
\bottomrule
\end{tabularx}
\end{center}


% ============================================================
% ENTREGABLES
% ============================================================
\section{Entregables}

El repositorio contiene:

\begin{itemize}
    \item código fuente del dominio;
    \item métodos de pago;
    \item patrones de diseño;
    \item motor de procesamiento;
    \item pruebas unitarias e integración;
    \item documentación de las ocho lecciones;
    \item spikes técnicos;
    \item diagrama de clases UML;
    \item README;
    \item LICENSE;
    \item configuración de pre-commit;
    \item workflow de integración continua;
    \item artefactos de distribución de la versión 1.0.0.
\end{itemize}

% ============================================================
% CONCLUSIÓN
% ============================================================
\section{Conclusión}

OrbitPay OO integra los conceptos fundamentales de Programación
Orientada a Objetos en una solución de dominio coherente,
extensible y verificable.

El proyecto evoluciona desde una problemática monolítica hacia
una arquitectura basada en abstracciones, encapsulamiento,
polimorfismo, composición y patrones de diseño.

La combinación de pruebas automatizadas, cobertura de código,
análisis estático, integración continua y empaquetado reproducible
proporciona evidencia técnica del funcionamiento de la solución
y de su preparación para distribución como paquete Python.

El release \textbf{OrbitPay OO 1.0.0} representa el cierre técnico
del proyecto integrador y reúne en un único artefacto las
competencias desarrolladas durante el módulo.

% ============================================================
% CIERRE
% ============================================================
\vfill

\begin{center}
    \rule{0.7\textwidth}{0.5pt}

    \vspace{0.3cm}

    \textbf{OrbitPay OO -- Release 1.0.0}

    \smallskip

    Luis Angel Leyva Castillo
\end{center}

\end{document}